% Source: https://artofproblemsolving.com/community/c3h1643422_fish_tank
% Stable identity: aops-topic-1643422-q01
Let the first tank have length $l$, width $w$ and height $h$, so its volume is $lwh$. The new tank has dimensions $kl$, $kw$, $kh$ and volume $(kl)(kw)(kh)=k^3\,lwh$.

\begin{asy}
usepackage("amsmath");
size(8cm);
defaultpen(fontsize(10pt));
pair ob=(0.45,0.3);
void box(pair o,real l,real h,real d,pen p){
 pair A=o,B=o+(l,0),C=o+(l,h),D=o+(0,h),v=d*ob;
 draw(A--B--C--D--cycle,p);draw(B--B+v--C+v--C,p);draw(D--D+v--C+v,p);draw(A+v--B+v^^A+v--D+v^^A--A+v,p+dashed);}
real k=0.25^(1/3);
box((0,0),4,2.4,2.5,linewidth(0.9));box((6.8,0),4*k,2.4*k,2.5*k,linewidth(0.9));
label("$l$",(2,0),S);label("$h$",(0,1.2),W);label("$w$",(4,0)+0.5*2.5*ob,SE);
label("$kl$",(6.8+2*k,0),S);label("$kh$",(6.8,1.2*k),W);label("$kw$",(6.8+4*k,0)+0.5*2.5*k*ob,SE);
label("$10$ gal",(2,-0.9));label("$2.5$ gal",(6.8+2*k,-0.9));
label("\mbox{}",(min(currentpicture,user=true)+max(currentpicture,user=true))/2,black); // keeps the text after the figure black when LaTeX typesets the labels inline
\end{asy}

The new volume is $\dfrac{1}{4}$ of the old, so
\[
k^3=\dfrac{1}{4},\qquad k=\sqrt[3]{\dfrac{1}{4}}=\dfrac{\sqrt[3]{2}}{2}\approx0.63 .
\]
So every side of the new tank is $\boxed{\dfrac{\sqrt[3]{2}}{2}\approx0.63}$ times the corresponding side of the old tank (not $\dfrac{1}{4}$ of it: shrinking each side to a quarter would give $\dfrac{1}{64}$ of the volume).

The actual side lengths cannot be found from the volume alone: many different boxes hold $10$ gallons (for instance long and shallow or tall and narrow). Only the scale factor is determined; once the old tank's dimensions are known, multiply each of them by $0.63$.
